\documentclass[pdflatex,sn-mathphys-num]{sn-jnl}

\usepackage{graphicx}
\usepackage{booktabs}
\usepackage{amsmath}
\usepackage{url}

\begin{document}

\title[Layer-Adaptive XMSSMT Parameterization]{Layer-Adaptive XMSSMT Parameterization: A Multi-Objective Study of WOTS Placement, Tree-Height Allocation, and Serialized State}

\author[1]{\fnm{Sohaib} \sur{Rana}}
\email{sohaibrana2023@student.usm.my}

\author*[1,2]{\fnm{Mohd Najwadi} \sur{Yusoff}}
\email{najwadi@usm.my}

\author*[3,4]{\fnm{Je Sen} \sur{Teh}}
\email{j.teh@deakin.edu.au}

\author[3]{\fnm{Adnan} \sur{Anwar}}
\email{adnan.anwar@deakin.edu.au}

\affil[1]{%
\orgdiv{School of Computer Sciences},
\orgname{Universiti Sains Malaysia},
\orgaddress{\city{Gelugor}, \state{Penang}, \country{Malaysia}}}

\affil[2]{%
\orgdiv{Cybersecurity Research Centre (CYRES)},
\orgname{Universiti Sains Malaysia},
\orgaddress{\city{Gelugor}, \state{Penang}, \country{Malaysia}}}

\affil[3]{%
\orgdiv{School of Information Technology},
\orgname{Deakin University},
\orgaddress{\city{Waurn Ponds}, \state{Victoria}, \country{Australia}}}

\affil[4]{%
\orgdiv{Deakin Cyber},
\orgname{Deakin University},
\orgaddress{\city{Waurn Ponds}, \state{Victoria}, \country{Australia}}}

\abstract{XMSSMT implementations commonly use uniform per-layer parameters, although lower and upper layers perform asymmetric computational and state-management roles. This study evaluates ordered layer-specific tree-height and WOTS parameterization in a two-layer XMSSMT FAST/BDS implementation with total height $H=20$. Five height allocations, $(6,14)$, $(8,12)$, $(10,10)$, $(12,8)$, and $(14,6)$, are combined with four ordered WOTS placements, $(16,16)$, $(16,4)$, $(4,16)$, and $(4,4)$, yielding 20 experimental configurations. A unified Raspberry Pi~5 protocol measures one complete lower-layer traversal cycle plus rollover for every configuration, producing 87,316 signing observations with matched verification, hardware-cycle, and logical-work measurements. Lower-layer WOTS placement dominates routine signing cost: across height allocations, $(4,16)$ reduces mean signing latency substantially relative to $(16,4)$ despite both producing 7075-byte signatures. Serialized-state size follows the exact relation $|\mathrm{SK}|=2654+120h_0+W_1$, while signature size follows $|\sigma|=675+W_0+W_1$; both models reproduce all 20 measured configurations exactly. Rollover behavior is composition-dependent rather than a fixed positive surcharge, with $(4,16)$ showing positive rollover increases and $(16,4)$ negative changes across all tested height allocations. An aggregate workload model explains measured CPU cycles with $R^2=0.9880$. Under the primary four-objective analysis of mean signing latency, P95 signing latency, signature size, and serialized state, three configurations form the strict Pareto front; two remain robust under 2\%, 5\%, and 10\% tolerance. The results quantify ordered layer effects in a concrete implementation and do not claim that heterogeneous XMSSMT parameterization itself is new.}

\keywords{XMSSMT, XMSS, hash-based signatures, post-quantum cryptography, WOTS+, Winternitz parameter, BDS traversal, heterogeneous parameterization, Raspberry Pi 5, multi-objective optimization}

\maketitle

\section{Introduction}
\label{sec:introduction}

Hash-based signatures are among the most mature approaches to
post-quantum digital signatures because their security can be based on
well-studied properties of cryptographic hash functions rather than on
number-theoretic assumptions. XMSS was introduced as a forward-secure
stateful hash-based signature scheme and later extended through the
multi-tree construction XMSSMT
\cite{Buchmann2011XMSS,RFC8391}. XMSS and XMSSMT use WOTS+ as their
one-time signature component \cite{Huelsing2013WOTSPlus,RFC8391}.

The multi-tree construction reduces the cost associated with very large
single Merkle trees by dividing the total tree height across multiple
layers. This introduces additional design parameters, including the
number of layers, subtree heights, Winternitz parameters, and traversal
parameters. These choices jointly affect signing cost, verification
cost, signature size, traversal state, and the computational behavior
at subtree boundaries.

Stateful hash-based signatures are standardized for restricted
deployment scenarios. RFC~8391 specifies XMSS and XMSSMT
\cite{RFC8391}, while RFC~8554 specifies LMS and its hierarchical
extension HSS \cite{RFC8554}. NIST SP~800-208 profiles XMSS, XMSSMT,
LMS, and HSS for stateful hash-based signature applications
\cite{NIST800208}. NIST emphasizes that such schemes require careful
state management because reuse of one-time signing state can invalidate
their security guarantees.

The XMSSMT parameter-selection problem has been studied previously.
In particular, H{\"u}lsing, Rausch, and Buchmann considered optimization
of layer-dependent XMSSMT parameters, including subtree heights,
traversal parameters, and Winternitz parameters
\cite{OptimalParametersXMSSMT}. Consequently, the use of heterogeneous
or layer-specific XMSSMT parameters is not itself a new construction.

A remaining implementation question is how the \emph{ordered placement}
of these parameters manifests in a concrete FAST/BDS implementation
when configurations are compared under an identical measurement
protocol. Two parameterizations can have the same total tree height and
even the same transmitted signature size while assigning computational
work differently between the frequently executed lower layer and the
less frequently refreshed upper layer. In addition, BDS traversal makes
signing cost index-dependent, so comparisons based only on short
early-index signing windows can obscure traversal bursts and rollover
behavior.

This work therefore studies a controlled two-layer XMSSMT design space
with fixed total height $H=20$ and layer-specific height and WOTS
parameters. Five ordered height allocations,
\[
(6,14),(8,12),(10,10),(12,8),(14,6),
\]
are combined with four ordered WOTS placements,
\[
(16,16),(16,4),(4,16),(4,4),
\]
to produce 20 experimental configurations. These configurations are
used as an implementation study and are not claimed to be standardized
XMSSMT parameter sets.

The evaluation differs from a parameter-only comparison in five ways.
First, all configurations are measured using one matched full-cycle
protocol covering a complete lower-layer traversal cycle and its
rollover boundary. Second, the implementation records both wall-clock
latency and hardware CPU cycles at operation level. Third, algorithmic
work counters are collected to relate observed performance to logical
XMSSMT operations. Fourth, signature and serialized-state sizes are
derived analytically and checked against every measured configuration.
Fifth, deployment trade-offs are evaluated using strict and
tolerance-based multi-objective Pareto analysis.

The main contributions are therefore as follows:
\begin{itemize}
    \item a FAST/BDS implementation of ordered layer-specific XMSSMT
    height and WOTS parameterization under a single matched
    full-cycle benchmark protocol;

    \item an operation-level characterization of lower-layer placement,
    traversal burstiness, and rollover behavior using latency, hardware
    cycles, and logical-work counters;

    \item exact analytical models for signature size and serialized
    implementation state, validated with zero-byte residuals across the
    complete 20-configuration matrix;

    \item an empirical aggregate workload model relating logical hash
    work and tree-node generation to measured processor cycles; and

    \item a multi-objective analysis of signing latency, tail latency,
    signature size, serialized state, and verification sensitivity,
    including robustness under symmetric $\epsilon$-dominance.
\end{itemize}

The contribution is thus not a new XMSSMT cryptographic primitive or the
first proposal of heterogeneous XMSSMT parameters. Rather, it is an
implementation-level quantification and validation of ordered
layer-placement effects, full-cycle traversal behavior, rollover
composition, serialized-state structure, and deployment trade-offs on
a Raspberry Pi~5 platform.


\section{Related Work}
\label{sec:related_work}

\subsection{XMSS, WOTS+, and XMSSMT}
\label{sec:related_xmss}

XMSS was introduced by Buchmann, Dahmen, and H{\"u}lsing as a practical
forward-secure hash-based signature scheme
\cite{Buchmann2011XMSS}. WOTS+ subsequently strengthened the
Winternitz one-time signature construction while reducing signature
length relative to earlier variants \cite{Huelsing2013WOTSPlus}.
RFC~8391 later specified interoperable WOTS+, XMSS, and XMSSMT
constructions \cite{RFC8391}.

XMSSMT organizes multiple XMSS trees hierarchically. A lower-layer tree
signs messages, while its root is authenticated by a WOTS+ signature in
the next layer. This organization changes how computation is distributed
across layers: lower-layer operations occur frequently, whereas
upper-layer authentication material can be reused until a lower subtree
is exhausted. The placement of layer-specific parameters can therefore
affect execution cost even when aggregate parameters such as total tree
height or total WOTS signature bytes remain unchanged.

\subsection{XMSSMT Parameter Optimization}
\label{sec:related_optimization}

H{\"u}lsing, Rausch, and Buchmann explicitly studied optimal parameter
selection for XMSSMT \cite{OptimalParametersXMSSMT}. Their formulation
allows parameters to vary across layers and considers the interaction
between subtree height, traversal parameter, and Winternitz parameter.
That work establishes important prior art for heterogeneous XMSSMT
parameterization and prevents interpreting layer-specific
$(h_i,k_i,w_i)$ choices as a novel concept in themselves.

The focus of the present work is different. Rather than introducing the
idea of heterogeneous parameters, we experimentally isolate the effect
of their \emph{ordered placement} in a concrete FAST/BDS implementation.
The study additionally measures an entire lower-layer traversal cycle,
including rollover, and records operation-level hardware cycles and
algorithmic work counters. This provides an empirical view of effects
that are structural in the scheme but can be obscured by limited
signing windows or by aggregate timing alone.

\subsection{Standardized Stateful Hash-Based Signatures}
\label{sec:related_standards}

RFC~8391 specifies XMSS and XMSSMT and defines interoperable parameter
sets for those schemes \cite{RFC8391}. The heterogeneous configurations
used in this work are experimental implementation variants and are not
assigned or represented as standardized RFC parameter identifiers.

LMS and HSS provide a second standardized stateful hash-based signature
family. RFC~8554 specifies LMS and the hierarchical HSS construction
\cite{RFC8554}. HSS is relevant to the present discussion because
hierarchical hash-based signatures generally expose design trade-offs
between lower-level signing work, upper-level authentication, key
lifetime, and state management.

NIST SP~800-208 profiles both families: XMSS/XMSSMT and LMS/HSS
\cite{NIST800208}. The recommendation approves selected parameter
sets rather than arbitrary experimental combinations and emphasizes the
operational importance of maintaining state correctly. Accordingly, the
parameter sweep in this study should be viewed as an implementation and
design-space experiment rather than as a proposal to redefine the
standardized profiles.

\subsection{Positioning of the Present Study}
\label{sec:related_positioning}

Prior work establishes the underlying XMSS construction, WOTS+,
multi-tree organization, and optimization of XMSSMT parameters.
The present study addresses a narrower empirical question: how do
ordered layer-specific height and WOTS choices interact with FAST/BDS
traversal, serialized implementation state, full-cycle signing cost,
and rollover behavior when all configurations are measured using an
identical operation-level protocol?

This distinction determines the novelty claim used throughout the paper.
The contribution lies in matched implementation evidence and analytical
validation, not in claiming the first use of heterogeneous XMSSMT
parameters. In particular, the comparison between $(w_0,w_1)=(4,16)$
and $(16,4)$ demonstrates why ordered placement deserves explicit
measurement: both configurations have the same transmitted signature
size, yet assign work differently between the continuously active lower
layer and the periodically refreshed upper layer.

\section{Methods}

\subsection{Experimental Design}
\label{sec:experimental_design}

We evaluate a two-layer XMSSMT construction with total Merkle-tree
height $H=20$ and $d=2$ layers. Unlike standardized XMSSMT parameter
sets, which use a common subtree height and Winternitz parameter across
layers, the configurations considered here are experimental
heterogeneous parameterizations intended to isolate the effects of
ordered layer placement. The layer heights satisfy
$h_0+h_1=20$, with
\[
(h_0,h_1)\in
\{(6,14),(8,12),(10,10),(12,8),(14,6)\},
\]
and the layer-specific WOTS parameters are
\[
(w_0,w_1)\in
\{(16,16),(16,4),(4,16),(4,4)\}.
\]
Their Cartesian product yields 20 configurations. No standardized OIDs
are claimed for these heterogeneous configurations.

\subsection{FAST/BDS Implementation}
\label{sec:fast_bds}

Measurements use the FAST XMSSMT implementation with BDS traversal.
For all evaluated configurations, the BDS parameter is $k=0$.
For $d=2$, the serialized traversal state contains the current
lower-layer state, the current upper-layer state, and a next
lower-layer state. The implementation additionally caches the
upper-layer WOTS signature used between lower-tree rollover events.
Consequently, serialized state size includes both traversal data and
cached public WOTS-signature material and should not be interpreted as
entirely secret memory.

The implementation invokes SHA-256 and SHA-512 through OpenSSL
\texttt{libcrypto.so.3}. The Raspberry Pi processor exposes the
AArch64 \texttt{sha2} feature. These observations establish the
software and processor capabilities used by the experiment but do not
by themselves identify the exact OpenSSL internal assembly path selected
at run time.

\subsection{Unified Full-Cycle Measurement Protocol}
\label{sec:fullcycle_protocol}

The revised evaluation uses a single matched measurement protocol for
all 20 configurations. This replaces earlier comparisons based on
different signing windows. For each configuration, ten key-generation
operations are measured. Signing is then measured over one complete
lower-layer traversal cycle together with the first signature after the
rollover boundary. Therefore, a configuration with lower-layer height
$h_0$ contributes
\[
2^{h_0}+1
\]
measured signing operations.

Each measured signature is immediately verified, so signing and
verification are evaluated over the same index sequence. The resulting
signing sample counts are 65, 257, 1025, 4097, and 16385 for
$h_0=6,8,10,12,$ and $14$, respectively. Across the complete
20-configuration matrix, the unified protocol produces 87,316 measured
signing operations.

Before measurement, three key-generation operations and 100 signing
operations are executed on a disposable key as an in-process warm-up.
The benchmark records the serialized signature index before and after
each signing operation and requires an increment of exactly one. This
provides an explicit state-progression check and identifies the
lower-layer rollover operation at
\[
i=2^{h_0}-1.
\]

\subsection{Platform and Timing}
\label{sec:platform}

Experiments were executed on a Raspberry Pi~5 Model B Rev.~1.1 with a
Broadcom BCM2712 processor using an AArch64 Cortex-A76 CPU. The system
ran Debian 13 (trixie) with the measured process pinned to CPU~3 under
the \texttt{performance} frequency governor. Observed CPU-frequency
samples were 2.4~GHz. Boost was disabled, and the platform throttling
indicator remained \texttt{0x0} throughout the campaign.

Wall-clock latency is measured with
\texttt{CLOCK\_MONOTONIC\_RAW}. Hardware CPU cycles are collected for
each key-generation, signing, and verification operation using
\texttt{perf\_event\_open()} with
\texttt{PERF\_COUNT\_HW\_CPU\_CYCLES}, excluding kernel and
hypervisor execution. The measured kernel remained
\texttt{6.18.34+rpt-rpi-2712} during the campaign.

\subsection{Logical Work Instrumentation}
\label{sec:logical_work}

In addition to latency and hardware cycles, the benchmark records
algorithm-level work counters for each operation: PRF invocations,
$F$-hash invocations, $H$-hash invocations, message-hash invocations,
WOTS chain steps, generated tree nodes, authentication-path nodes, and
BDS updates. We define
\[
N_{\mathrm{logical}}
 =
N_{\mathrm{PRF}}
+N_F
+N_H
+N_{\mathrm{msg}}.
\]
These counters denote logical XMSS/XMSSMT hash-function invocations and
must not be interpreted as counts of SHA-256 compression-function
calls.

The instrumentation shows that WOTS chain steps and $F$-hash calls are
identical in the present implementation. Strong collinearity also exists
among several primitive counters. Therefore, primitive-level regression
coefficients are not interpreted as independent operation costs.

\subsection{Summary Statistics}
\label{sec:statistics}

Because BDS traversal produces index-dependent and bursty signing
costs, signing observations are treated as an ordered traversal
sequence rather than as independent and identically distributed
samples. For each configuration we report the arithmetic mean, median,
sample standard deviation, coefficient of variation, minimum, maximum,
and 95th percentile. Percentiles use linear interpolation at position
$(n-1)p$.

The complete lower-layer cycle is retained without deleting
wall-clock outliers. Hardware-cycle measurements and logical work
counters are reported alongside latency to distinguish algorithmic
work variation from host-level scheduling or operating-system effects.

\subsection{Rollover Analysis}
\label{sec:rollover_analysis}

For every configuration, the rollover operation is defined as the
signature whose serialized index before signing equals
$2^{h_0}-1$. Its cost is compared with the mean of the immediately
preceding and following signing operations. We report
\[
\Delta_{\mathrm{roll}}
=
100\,
\frac{T_{\mathrm{roll}}-
      (T_{\mathrm{pre}}+T_{\mathrm{post}})/2}
     {(T_{\mathrm{pre}}+T_{\mathrm{post}})/2}.
\]

A rollover is not assumed to impose a fixed positive surcharge.
At the lower-tree boundary the implementation changes traversal-state
composition and generates a fresh upper-layer WOTS signature. The
resulting latency change can therefore be positive or negative,
depending on the ordered WOTS parameters and the work displaced at the
boundary.

\subsection{Analytical Artifact-Size Model}
\label{sec:size_model}

For $n=32$ and $k=0$, the serialized BDS state for the even subtree
heights used in this study reduces to
\[
S(h)=120h+41 \quad \text{bytes}.
\]

For the two-layer FAST topology, the serialized XMSSMT state is
\[
|\mathrm{SK}|
 =
131+2S(h_0)+S(h_1)+W_1,
\]
where $W_1$ is the cached upper-layer WOTS-signature size. Since
$h_0+h_1=20$,
\[
|\mathrm{SK}|=2654+120h_0+W_1.
\]

The WOTS signature sizes are 4256 bytes for $w=4$ and 2144 bytes for
$w=16$. Hence,
\[
|\mathrm{SK}|=
\begin{cases}
4798+120h_0, & w_1=16,\\
6910+120h_0, & w_1=4.
\end{cases}
\]

The XMSSMT signature size is
\[
|\sigma|
=
675+W_0+W_1,
\]
which yields 4963, 7075, 7075, and 9187 bytes for
$(w_0,w_1)=(16,16),(16,4),(4,16),$ and $(4,4)$, respectively.
The analytical signature and serialized-state formulas reproduce all
20 measured artifact sizes exactly, with zero-byte residuals.

\subsection{Aggregate Workload Model}
\label{sec:aggregate_model}

To relate logical algorithmic work to measured processor cost, an
aggregate empirical model is fitted over all 87,316 signing
observations:
\[
C =
\alpha
+\beta N_{\mathrm{logical}}
+\gamma N_{\mathrm{tree}},
\]
where $C$ is measured CPU cycles and $N_{\mathrm{tree}}$ is the number
of generated tree nodes.

For the Raspberry Pi~5 campaign, the fitted model is
\[
C =
-157535.67
+1779.238\,N_{\mathrm{logical}}
+56185.960\,N_{\mathrm{tree}}.
\]
The model achieves $R^2=0.9880$, a median absolute percentage error of
4.348\%, and a 95th-percentile absolute percentage error of 6.757\%.
These coefficients are platform- and workload-specific empirical
calibration parameters and are not interpreted as universal costs of
the underlying hash primitives.

\subsection{Multi-Objective Analysis}
\label{sec:pareto_method}

The primary deployment-oriented Pareto analysis minimizes four
objectives simultaneously: mean signing latency, P95 signing latency,
signature size, and serialized state size. A strict Pareto front is
computed first. Robustness is then examined using symmetric
$\epsilon$-dominance for $\epsilon\in\{2\%,5\%,10\%\}$.

Under this rule, configuration $A$ is considered no worse than
configuration $B$ when every objective of $A$ is at most
$(1+\epsilon)$ times the corresponding objective of $B$, with at least
one objective smaller than $(1-\epsilon)$ times that of $B$. A
secondary sensitivity analysis additionally includes mean verification
latency as a fifth objective.

\begin{table*}[t]
\centering
\caption{Full-cycle key-generation, signing, verification, rollover, and artifact-size results for the 20 heterogeneous XMSSMT configurations on Raspberry Pi~5. KeyGen reports mean latency $\pm$ sample SD in seconds ($n=10$ for the original campaign cells; $n=30$ for the pooled $h_0=14$ randomized validation cells). Signing mean, median, and P95 are computed over one complete lower-layer traversal cycle plus the first signature after rollover. For $h_0=14$, signing and verification statistics are pooled over three independently randomized full-cycle repetitions ($n=49{,}155$ per configuration), and Rollover $\Delta$ is the mean of the three neighbour-relative repeat deltas; the original single-run values are reported in Table~\ref{tab:h14_randomized_validation}. Serialized state includes traversal state and cached upper-layer WOTS signatures.}
\label{tab:main_20cell_results}
\scriptsize
\setlength{\tabcolsep}{1.65pt}
\begin{tabular}{@{}rrrrrrrrrrrrr@{}}
\toprule
$h_0$ & $h_1$ & $w_0$ & $w_1$ & KeyGen & \multicolumn{3}{c}{Sign latency (ms)} & Cycles & Verify & Rollover & Sig. & State \\
\cmidrule(lr){6-8}
& & & & (s) & Mean & Median & P95 & (M) & (ms) & $\Delta$ (\%) & (B) & (B) \\
\midrule
6 & 14 & 16 & 16 & $39.890\pm0.038$ & 8.482 & 7.064 & 13.266 & 20.339 & 2.607 & 18.33 & 4963 & 5518 \\
6 & 14 & 16 & 4 & $22.102\pm0.010$ & 8.346 & 6.134 & 13.107 & 20.018 & 2.161 & -8.45 & 7075 & 7630 \\
6 & 14 & 4 & 16 & $39.315\pm0.022$ & 4.623 & 4.528 & 7.233 & 11.087 & 2.242 & 68.99 & 7075 & 5518 \\
6 & 14 & 4 & 4 & $22.097\pm0.005$ & 4.599 & 3.659 & 7.247 & 11.030 & 1.608 & 13.57 & 9187 & 7630 \\
\addlinespace
8 & 12 & 16 & 16 & $10.429\pm0.007$ & 10.674 & 12.957 & 15.469 & 25.598 & 2.549 & 20.44 & 4963 & 5758 \\
8 & 12 & 16 & 4 & $6.099\pm0.003$ & 10.675 & 12.988 & 15.496 & 25.601 & 2.123 & -9.33 & 7075 & 7870 \\
8 & 12 & 4 & 16 & $10.259\pm0.002$ & 5.966 & 7.278 & 8.675 & 14.308 & 2.146 & 71.62 & 7075 & 5758 \\
8 & 12 & 4 & 4 & $5.841\pm0.003$ & 5.892 & 7.189 & 8.581 & 14.129 & 1.707 & 17.83 & 9187 & 7870 \\
\addlinespace
10 & 10 & 16 & 16 & $4.888\pm0.001$ & 13.066 & 15.435 & 17.913 & 31.336 & 2.605 & 18.64 & 4963 & 5998 \\
10 & 10 & 16 & 4 & $3.816\pm0.001$ & 13.030 & 15.389 & 17.874 & 31.249 & 2.168 & -6.57 & 7075 & 8110 \\
10 & 10 & 4 & 16 & $3.818\pm0.001$ & 7.240 & 8.571 & 9.947 & 17.364 & 2.120 & 71.36 & 7075 & 5998 \\
10 & 10 & 4 & 4 & $2.746\pm0.003$ & 7.227 & 8.555 & 9.932 & 17.332 & 1.674 & 16.97 & 9187 & 8110 \\
\addlinespace
12 & 8 & 16 & 16 & $10.399\pm0.008$ & 15.469 & 17.838 & 20.299 & 37.022 & 2.746 & 21.43 & 4963 & 6238 \\
12 & 8 & 16 & 4 & $10.148\pm0.008$ & 15.430 & 17.826 & 20.283 & 36.996 & 2.153 & -10.94 & 7075 & 8350 \\
12 & 8 & 4 & 16 & $6.090\pm0.002$ & 8.532 & 9.869 & 11.232 & 20.460 & 2.177 & 64.13 & 7075 & 6238 \\
12 & 8 & 4 & 4 & $5.822\pm0.003$ & 8.541 & 9.882 & 11.241 & 20.482 & 1.649 & 17.54 & 9187 & 8350 \\
\addlinespace
14 & 6 & 16 & 16 & $39.277\pm0.053$ & 17.823 & 20.250 & 22.686 & 42.737 & 2.587 & 19.78 & 4963 & 6478 \\
14 & 6 & 16 & 4 & $39.373\pm0.209$ & 17.889 & 20.326 & 22.883 & 42.896 & 2.157 & -8.99 & 7075 & 8590 \\
14 & 6 & 4 & 16 & $22.127\pm0.019$ & 9.907 & 11.264 & 12.620 & 23.757 & 2.110 & 70.89 & 7075 & 6478 \\
14 & 6 & 4 & 4 & $22.103\pm0.078$ & 9.954 & 11.296 & 12.686 & 23.867 & 1.672 & 18.61 & 9187 & 8590 \\
\bottomrule
\end{tabular}
\end{table*}


\begin{table*}[t]
\centering
\caption{Pareto-front size under strict, relaxed, and margin-based
dominance sensitivity for the three deployment objective scenarios.
Scenario~A minimizes mean signing latency, P95 signing latency,
signature size, and serialized state size; Scenario~B additionally
includes mean key-generation latency; Scenario~C additionally includes
mean verification latency. For relaxed dominance, a candidate may be
within $(1+\epsilon)$ of another configuration in every objective while
improving at least one objective by more than $(1-\epsilon)$. Margin
dominance instead requires an improvement of at least $\epsilon$ in
every included objective. All objectives are minimized.}
\label{tab:pareto_robustness}
\footnotesize
\setlength{\tabcolsep}{3.2pt}
\begin{tabular}{@{}lrrrrrrr@{}}
\toprule
& & \multicolumn{3}{c}{Relaxed dominance}
& \multicolumn{3}{c}{Margin dominance} \\
\cmidrule(lr){3-5}
\cmidrule(lr){6-8}
Scenario
& Strict
& 2\% & 5\% & 10\%
& 2\% & 5\% & 10\% \\
\midrule
A: Sign/P95/Sig./State
& 3 & 2 & 2 & 2 & 12 & 12 & 12 \\
B: A + Key generation
& 10 & 9 & 9 & 9 & 16 & 16 & 16 \\
C: B + Verification
& 12 & 9 & 9 & 9 & 20 & 20 & 20 \\
\bottomrule
\end{tabular}
\end{table*}


\section{Results}
\label{sec:results}

\subsection{Full-Cycle Signing Cost and Layer Placement}
\label{sec:results_fullcycle}

Table~\ref{tab:main_20cell_results} summarizes the matched full-cycle
measurements for all 20 heterogeneous XMSSMT configurations.
Figure~\ref{fig:sign_h0} shows mean signing latency as a function of the
lower-layer height $h_0$ for each ordered WOTS placement.

Across all four WOTS placements, increasing $h_0$ generally increases
full-cycle signing cost. For $(w_0,w_1)=(16,16)$, mean signing latency
rises from 8.482~ms at $(h_0,h_1)=(6,14)$ to 17.891~ms at
$(14,6)$. The corresponding mean hardware-cycle count rises from
20.339 to 42.907 million cycles. The same structural trend appears for
the $w_0=4$ configurations, although at a substantially lower signing
cost.

The lower-layer WOTS parameter has the dominant effect on ordinary
full-cycle signing. With $w_1=16$, changing $w_0$ from 16 to 4 reduces
mean signing latency from 8.482 to 4.623~ms at $h_0=6$, from 10.674 to
5.966~ms at $h_0=8$, from 13.066 to 7.240~ms at $h_0=10$, from
15.469 to 8.532~ms at $h_0=12$, and from 17.891 to 9.907~ms at
$h_0=14$. These reductions are approximately 44--46\% across the tested
height allocations.

In contrast, changing only $w_1$ generally has a much smaller effect on
the full-cycle mean because the upper-layer WOTS signature is cached
between lower-tree rollovers rather than regenerated for every
signature. This distinction is visible, for example, at
$(h_0,h_1)=(10,10)$: mean signing latency is 13.066~ms for
$(16,16)$ and 13.030~ms for $(16,4)$, while it is 7.240~ms for
$(4,16)$ and 7.227~ms for $(4,4)$.

The median and P95 measurements also reveal the index-dependent
burstiness of BDS traversal. Consequently, the mean alone does not
fully characterize signing behavior. For example, the
$(6,14,16,16)$ configuration has a mean of 8.482~ms, a median of
7.064~ms, and a P95 of 13.266~ms. These differences motivate retaining
median and tail statistics alongside mean latency.

\begin{figure*}[t]
\centering
\includegraphics[width=0.78\textwidth]
{FIG_SIGN_MEAN_VS_H0_BY_WOTS.pdf}
\caption{Mean full-cycle signing latency versus lower-layer height
$h_0$ for the four ordered WOTS placements. Each point is computed from
one complete lower-layer traversal cycle plus the first signature after
rollover.}
\label{fig:sign_h0}
\end{figure*}


\subsection{Exact Artifact-Size Structure}
\label{sec:results_sizes}

The measured signature and serialized-state sizes follow the analytical
models exactly. For fixed total height $H=20$, signature size depends
on the ordered WOTS parameters but not on how the total tree height is
partitioned between the two layers:
\[
|\sigma|=675+W_0+W_1.
\]
This gives 4963~bytes for $(16,16)$, 7075~bytes for both $(16,4)$ and
$(4,16)$, and 9187~bytes for $(4,4)$.

The equality of the 7075-byte signatures for $(16,4)$ and $(4,16)$ is
particularly useful for separating communication cost from computation.
Although the two placements transmit identical signature sizes, their
signing behavior differs substantially because the lower-layer WOTS
parameter determines work performed for every signature.

Serialized state follows
\[
|\mathrm{SK}|=2654+120h_0+W_1.
\]
Thus, increasing $h_0$ by one increases serialized state by exactly
120~bytes for fixed $w_1$. Changing the upper-layer parameter from
$w_1=16$ to $w_1=4$ increases state by exactly 2112~bytes, corresponding
to the difference between the cached upper-layer WOTS-signature sizes.

Across all 20 configurations, both analytical formulas reproduce the
measured sizes with zero-byte residuals. The serialized-state quantity
includes traversal state and cached upper-layer WOTS-signature material;
it is therefore an implementation-state measure rather than a measure
of exclusively secret data.


\subsection{Rollover Is a Work-Composition Effect}
\label{sec:results_rollover}

Figure~\ref{fig:rollover} reports the latency change at the lower-tree
rollover boundary relative to the mean of the immediately preceding and
following signatures. The results show that rollover is not a fixed
positive surcharge.

For $(w_0,w_1)=(4,16)$, rollover latency increases by
64.13--74.29\% across the tested height allocations. By contrast,
$(16,4)$ produces negative rollover changes for every height allocation,
ranging from $-10.94$\% to $-6.57$\%. The $(16,16)$ configurations
show positive changes of 18.33--24.04\%, while $(4,4)$ produces
smaller positive changes of 13.57--17.83\%.

The hardware-cycle and logical-work measurements follow the same
qualitative pattern. Averaged across height allocations, the
$(4,16)$ placement increases rollover latency by approximately
70.08\%, hardware cycles by 70.09\%, and logical hash work by 71.28\%.
For $(16,4)$, the corresponding averages are approximately
$-8.53$\%, $-8.53$\%, and $-8.92$\%, respectively. This agreement
indicates that the placement-dependent rollover effect is primarily
structural rather than an artifact of wall-clock noise.

The counter decomposition further shows that rollover changes the mix
of lower-layer traversal work and upper-layer WOTS work. Because the
boundary operation both changes BDS state and generates a fresh
upper-layer WOTS signature, its cost cannot be represented accurately
as a single constant additive penalty.

\begin{figure*}[t]
\centering
\includegraphics[width=0.78\textwidth]
{FIG_ROLLOVER_DELTA_VS_H0_BY_WOTS.pdf}
\caption{Rollover latency change relative to the mean of the immediately
preceding and following signatures. Negative values indicate that the
rollover signature is faster than its local adjacent-signature
baseline.}
\label{fig:rollover}
\end{figure*}


\subsection{Aggregate Workload Model}
\label{sec:results_model}

The primitive instrumentation counters are strongly collinear.
In particular, WOTS chain steps are identical to $F$-hash calls in the
instrumented implementation, while PRF, $F$-hash, and $H$-hash counts
also exhibit substantial dependence. Primitive-level regression
coefficients are therefore not interpreted as independent operation
costs.

Instead, the aggregate model
\[
C=
\alpha+
\beta N_{\mathrm{logical}}+
\gamma N_{\mathrm{tree}}
\]
is fitted over all 87,316 signing observations. The resulting model is
\[
C=
-157535.67
+1779.238\,N_{\mathrm{logical}}
+56185.960\,N_{\mathrm{tree}}.
\]

The fit achieves
\[
R^2=0.9880,
\]
with a median absolute percentage error of 4.348\% and a
95th-percentile absolute percentage error of 6.757\%.
Figure~\ref{fig:model_validation} compares measured and predicted CPU
cycles.

These coefficients should be interpreted as an empirical calibration
of this Raspberry Pi~5 implementation and workload. In particular,
the coefficient associated with logical hash invocations is not a
measurement of the universal cost of one SHA-256 compression or one
cryptographic hash primitive.

\begin{figure*}[t]
\centering
\includegraphics[width=0.72\textwidth]
{FIG_AGGREGATE_MODEL_MEASURED_VS_PREDICTED.pdf}
\caption{Measured versus predicted signing CPU cycles for the aggregate
workload model over all 87,316 signing observations. The dashed line
denotes ideal prediction.}
\label{fig:model_validation}
\end{figure*}


\subsection{Multi-Objective Deployment Trade-offs}
\label{sec:results_pareto}

The primary Pareto analysis jointly minimizes mean signing latency,
P95 signing latency, signature size, and serialized state size.
Three configurations lie on the strict four-objective front:
$(6,14,16,16)$, $(6,14,4,16)$, and $(6,14,4,4)$.

These configurations represent different trade-offs rather than a
single universally preferable parameterization. The
$(6,14,16,16)$ configuration has the smallest signature among the
three at 4963~bytes, with a mean signing latency of 8.482~ms and
serialized state of 5518~bytes. The $(6,14,4,16)$ configuration reduces
mean signing latency to 4.623~ms while retaining the same 5518-byte
serialized state, at the cost of increasing the signature to
7075~bytes. The $(6,14,4,4)$ configuration has a slightly lower mean
signing latency of 4.599~ms and a mean verification latency of
1.608~ms, but increases the signature to 9187~bytes and serialized
state to 7630~bytes.

Tolerance analysis shows that the strict front is partly sensitive to
small differences. Under symmetric 2\%, 5\%, and 10\%
$\epsilon$-dominance, $(6,14,16,16)$ and $(6,14,4,16)$ remain on the
front at every tested tolerance, whereas $(6,14,4,4)$ is retained only
under strict dominance. Table~\ref{tab:pareto_robustness} reports this
sensitivity.

Figure~\ref{fig:pareto} visualizes the relation between signing latency,
signature size, serialized state, and strict Pareto membership. The
results reinforce that WOTS placement cannot be selected from signature
size alone: $(16,4)$ and $(4,16)$ both produce 7075-byte signatures,
yet their signing and rollover costs differ markedly.

A secondary analysis includes mean verification latency as a fifth
objective. The strict front expands because small verification
differences prevent several configurations from dominating one another.
However, under 5\% and 10\% tolerance the verification-inclusive front
contracts to the same three $(6,14)$ configurations identified above.
This sensitivity indicates that conclusions based solely on strict
dominance should be interpreted cautiously when objective differences
are small.

\begin{figure*}[t]
\centering
\includegraphics[width=0.72\textwidth]
{FIG_PARETO_LATENCY_SIGNATURE_STATE.pdf}
\caption{Deployment trade-off between mean signing latency and signature
size. Marker area reflects serialized-state size; strict Pareto
configurations are labelled.}
\label{fig:pareto}
\end{figure*}


\subsection{Timing Outlier and Robustness of Interpretation}
\label{sec:results_outlier}

The largest observed wall-clock signing latency is 45.239~ms for
$(h_0,h_1,w_0,w_1)=(12,8,16,16)$. This observation is retained in the
dataset. It occurs within a short cluster of elevated wall-clock
measurements around the same signature indices.

The corresponding cycle increase is substantially smaller than the
wall-clock increase: the maximum latency is approximately 2.54 times
the local median, whereas the hardware-cycle ratio is approximately
1.17. During the campaign the measured temperature remained within the
observed operating range and the throttling indicator remained
\texttt{0x0}. The available measurements are therefore consistent with
a host-level scheduling, interrupt, or operating-system disturbance,
although the experiment does not record sufficient information to
identify a specific cause.

For this reason, the observation is not deleted. Median, P95, hardware
cycles, and logical-work counters are reported alongside mean and
maximum latency so that deployment conclusions do not depend on a
single wall-clock extreme.

\section{Discussion}
\label{sec:discussion}

\subsection{Implications of Ordered Layer Parameterization}
\label{sec:discussion_ordering}

The results show that heterogeneous XMSSMT parameterization cannot be
characterized solely by the unordered set of layer parameters.
Placement matters because the two layers perform different roles during
normal signing and at lower-tree rollover boundaries.

The strongest example is the comparison between
$(w_0,w_1)=(4,16)$ and $(16,4)$. Both placements produce the same
7075-byte signature because the transmitted signature contains one WOTS
signature from each layer and therefore depends on the sum
$W_0+W_1$. Their execution costs, however, differ substantially.
The lower-layer WOTS signature is generated for every signing operation,
whereas the upper-layer WOTS signature is cached and refreshed only at
a lower-tree rollover. Consequently, placing $w=4$ at the lower layer
reduces routine signing cost, while placing it only at the upper layer
does not provide the same full-cycle benefit.

This distinction is deployment-relevant because two configurations with
the same communication cost can have substantially different signing
latency, processor-cycle demand, serialized state, and rollover
behavior. Signature size alone is therefore insufficient for selecting
an ordered heterogeneous XMSSMT configuration.

\subsection{Why Lower-Layer Height Dominates Full-Cycle Cost}
\label{sec:discussion_height}

Increasing $h_0$ raises the amount of lower-layer traversal work and
lengthens the interval between upper-layer rollover events. In the
evaluated FAST/BDS implementation with $k=0$, the initial TreeHash
update budget is
\[
U_0=\frac{h_0}{2}.
\]
The serialized traversal topology also stores the lower-layer BDS state
twice---as current and next state---but the upper-layer state only once.
This asymmetry explains why lower-layer height affects both computation
and serialized-state size more strongly than an equivalent increase in
upper-layer height.

For fixed total height $H=20$, increasing $h_0$ by one increases the
derived serialized state by exactly 120 bytes. This relation is not an
empirical regression: it follows directly from the implemented BDS
serialization structure for $n=32$ and $k=0$ and is reproduced exactly
by all 20 measured configurations.

The performance trend is similarly consistent across the parameter
matrix. Moving from $(h_0,h_1)=(6,14)$ toward $(14,6)$ generally
increases mean signing latency and hardware cycles. Thus, for this
implementation and workload, allocating a larger fraction of the fixed
total height to the lower layer trades less frequent rollover for higher
per-signature traversal cost and larger serialized state.

\subsection{Rollover Cost Is Not a Constant Penalty}
\label{sec:discussion_rollover}

A key consequence of the full-cycle protocol is that rollover behavior
can be observed directly rather than inferred from early signature
indices. The measured rollover change is strongly dependent on ordered
WOTS placement and is not uniformly positive.

For $(4,16)$, the rollover signature is substantially more expensive
than its adjacent-signature baseline. For $(16,4)$, the rollover
signature is consistently faster than the local baseline across all
five tested height allocations. The corresponding hardware-cycle and
logical-work changes closely track the wall-clock changes.

This behavior follows from the work composition at the boundary.
Rollover does not simply add one upper-layer WOTS computation to an
otherwise identical signing operation. It also changes the BDS state,
swaps the precomputed next lower-layer state into the current position,
reinitializes traversal work, and changes the mixture of lower- and
upper-layer operations. Accordingly, the rollover effect should be
described as a boundary-specific composition of work rather than as a
fixed surcharge.

This interpretation also explains why immediate-neighbor comparison is
more informative than quoting rollover latency in isolation. The local
baseline partially controls for the index-dependent burst structure of
BDS traversal.

\subsection{Relationship Between Logical Work and Processor Cost}
\label{sec:discussion_model}

The aggregate workload model provides a compact explanation for most of
the measured cycle variation. Over 87,316 signing observations, logical
hash work and generated tree nodes explain approximately 98.8\% of the
variance in measured CPU cycles.

The model is useful as an implementation-level validation tool rather
than as a primitive benchmark. The logical counters intentionally
represent XMSS/XMSSMT algorithmic operations. They do not count
SHA-256 compression calls, and several primitive counters are strongly
collinear. In particular, WOTS chain steps and $F$-hash calls are
identical in the instrumented code. Therefore, interpreting separate
primitive regression coefficients as independent cryptographic
operation costs would be misleading.

The aggregate model avoids that interpretation and instead establishes
that the measured performance differences are closely associated with
changes in algorithmic work. The remaining prediction error includes
implementation overhead, instruction-level effects, memory behavior,
operating-system activity, and other costs not represented by the two
aggregate predictors.

\subsection{Deployment Trade-offs}
\label{sec:discussion_deployment}

The Pareto analysis illustrates that heterogeneous parameter selection
is inherently multi-objective. No single configuration simultaneously
minimizes signing latency, verification latency, signature size, and
serialized state.

The strict four-objective front contains three configurations from the
$(h_0,h_1)=(6,14)$ family. The $(16,16)$ WOTS placement provides the
smallest signature among them, whereas $(4,16)$ substantially reduces
signing latency while preserving the same serialized-state size.
The $(4,4)$ placement slightly reduces mean signing latency further but
requires both a larger signature and larger serialized state.

The tolerance analysis is important because strict Pareto membership can
be unstable when objective differences are small. The
$(6,14,4,4)$ configuration disappears from the front under 2\%,
5\%, and 10\% symmetric tolerance, whereas $(6,14,16,16)$ and
$(6,14,4,16)$ remain. This indicates that the latter two configurations
represent more robust trade-off points within the measured objective
set.

The verification-inclusive analysis further demonstrates the sensitivity
of strict dominance to additional objectives. Adding mean verification
latency expands the strict front, but a 5\% or 10\% tolerance again
reduces the front to the same three $(6,14)$ configurations. These
results support reporting both strict and tolerance-based Pareto
analyses rather than drawing deployment conclusions from a single
strict front.

\subsection{Contribution Relative to Prior XMSSMT Parameter Studies}
\label{sec:discussion_novelty}

The contribution of this study is not the claim that layer-specific
XMSSMT parameters are themselves new. Earlier XMSSMT parameter work has
considered layer-dependent quantities, including subtree heights,
traversal parameters, and Winternitz parameters. The present study
instead focuses on an implementation and measurement gap.

Specifically, the study combines layer-specific height and WOTS
parameterization with a FAST/BDS XMSSMT implementation; evaluates
ordered placement effects under one matched full-cycle protocol;
measures both wall-clock latency and hardware CPU cycles at operation
level; records algorithmic work counters; validates serialized-state
and signature-size models exactly; measures lower-tree rollover
behavior directly; and evaluates deployment trade-offs using strict
and tolerance-based Pareto analysis.

This framing is important because several of the observed effects are
structural consequences of XMSSMT and WOTS organization rather than
new cryptographic mechanisms. The empirical contribution is the
quantification and validation of those effects under a controlled,
matched Raspberry Pi~5 implementation.

\subsection{Limitations}
\label{sec:limitations}

Several limitations constrain the scope of the conclusions.

First, the evaluated configurations are experimental heterogeneous
XMSSMT configurations. They are not standardized XMSSMT parameter sets,
and no standardized OIDs are claimed for them. The results therefore
characterize implementation and deployment trade-offs rather than
proposing direct replacements for standardized parameter sets.

Second, the study fixes $H=20$, $d=2$, $n=32$, and BDS $k=0$. Other
total heights, numbers of layers, digest sizes, or traversal parameters
may change both the absolute costs and the relative importance of
layer placement. In particular, increasing $d$ introduces additional
upper-layer states and rollover levels that are not represented by the
two-layer experiment.

Third, only $w\in\{4,16\}$ is used in the final 20-cell heterogeneous
matrix. These values are sufficient to create clear placement contrasts,
but they do not exhaust the possible Winternitz parameter space.

Fourth, measurements were obtained on one Raspberry Pi~5 platform using
one compiler and one software stack. The implementation invokes
OpenSSL SHA-256/SHA-512 through \texttt{libcrypto.so.3}, and the CPU
exposes the AArch64 \texttt{sha2} feature, but the experiment does not
trace the exact OpenSSL internal implementation path selected for each
hash invocation. Consequently, the empirical cycle-model coefficients
should not be generalized to other processors or cryptographic
libraries.

Fifth, BDS signing costs are index-dependent and serially structured.
The observations in a full traversal cycle are therefore not
independent and identically distributed samples. The study reports
descriptive statistics, hardware cycles, and logical work rather than
applying inferential procedures that assume independent signing
observations.

Sixth, one configuration contains a pronounced wall-clock timing
outlier. The observation is retained because the experiment provides no
evidence of invalid cryptographic execution or thermal throttling.
Hardware cycles and logical work indicate that its wall-clock magnitude
is not proportional to algorithmic work, making host-level scheduling
or interrupt activity a plausible explanation. However, the precise
cause was not instrumented and therefore cannot be established.

Finally, the Pareto results depend on the selected objective set.
Alternative deployment contexts may assign different importance to
key-generation latency, verification throughput, energy, RAM usage,
flash footprint, state-write durability, or communication bandwidth.
The reported Pareto sets should therefore be interpreted as
objective-dependent trade-off surfaces rather than universal parameter
recommendations.

\subsection{Summary of Design Implications}
\label{sec:discussion_summary}

Within the evaluated design space, four implementation-level conclusions
are supported consistently by the analytical and empirical results.

First, lower-layer WOTS placement governs routine signing cost more
strongly than upper-layer WOTS placement because the lower-layer WOTS
signature is regenerated on every signature.

Second, lower-layer height controls both traversal work and serialized
state through the asymmetric FAST/BDS state topology.

Third, rollover behavior depends on the composition of work at the
boundary and can be either more expensive or less expensive than nearby
signatures.

Fourth, communication cost, state cost, and signing cost cannot be
inferred from one another. Parameter selection therefore requires
joint consideration of ordered layer structure and deployment
objectives.

\section{Conclusion}
\label{sec:conclusion}

This study evaluated ordered layer-specific tree-height and WOTS
parameterization in a two-layer XMSSMT FAST/BDS implementation using a
matched full-cycle Raspberry Pi~5 benchmark. The revised protocol
measures one complete lower-layer traversal cycle plus rollover for
each of 20 experimental configurations, producing 87,316 signing
observations together with matched verification measurements,
hardware-cycle counts, and logical-work counters.

The results show that ordered placement matters. Lower-layer WOTS
selection dominates routine signing cost, whereas upper-layer WOTS
selection has a stronger influence on rollover behavior and serialized
state. Configurations with $(w_0,w_1)=(4,16)$ and $(16,4)$ produce the
same 7075-byte signature, yet their computational and state behavior
differ substantially because the two layers perform different roles.

Tree-height allocation also produces a systematic structural effect.
For fixed total height $H=20$, increasing the lower-layer height raises
serialized state by 120 bytes per level under the measured FAST/BDS
topology, while full-cycle signing cost increases because additional
lower-layer traversal work is exercised more frequently.

The analytical artifact-size model reproduces all measured signature
and serialized-state sizes exactly. The aggregate workload model also
explains most measured CPU-cycle variation, with $R^2=0.9880$, while
remaining explicitly an empirical platform/workload calibration rather
than a primitive hash-cost model.

Rollover measurements further show that subtree transition cost is not
a constant positive surcharge. Its effect depends on the ordered WOTS
placement and on how routine traversal work and upper-layer WOTS work
are composed at the boundary.

Under the primary four-objective analysis of mean signing latency,
P95 signing latency, signature size, and serialized state, three
configurations lie on the strict Pareto front. Two of these,
$(6,14,16,16)$ and $(6,14,4,16)$, remain on the front under 2\%, 5\%,
and 10\% symmetric tolerance, whereas $(6,14,4,4)$ is retained only
under strict dominance. These results reinforce that XMSSMT parameter
selection is deployment-dependent and should not be reduced to a single
universal optimum.

The contribution is therefore not a new XMSSMT primitive or the first
use of heterogeneous parameters. It is a matched implementation-level
quantification of ordered layer effects, traversal behavior, rollover
composition, serialized-state structure, and multi-objective deployment
trade-offs. Future work should extend the analysis to additional total
heights, layer counts, Winternitz parameters, BDS settings, processor
classes, and standardized deployment contexts.

\bmhead{Acknowledgements}
The authors acknowledge the support of the School of Computer Sciences, Universiti Sains Malaysia, and the Cybersecurity Research Centre (CYRES).

\section*{Declarations}

\bmhead{Author contributions}
Sohaib Rana: Conceptualization, methodology, software, validation, formal analysis, investigation, data curation, visualization, and writing -- original draft.

Mohd Najwadi Yusoff: Supervision, conceptualization, methodology, project administration, and writing -- review and editing.

Je Sen Teh: Supervision, methodology, validation, and writing -- review and editing.

Adnan Anwar: Validation and writing -- review and editing.

\bmhead{Funding}
The authors received no specific funding for this work.

\bmhead{Data availability}
The processed benchmark datasets, normalized operation-level measurements, analysis outputs, provenance records, and integrity manifests supporting this study are provided in Online Resource 1 and in the public reproducibility repository at \url{https://github.com/sohaibrana1/xmssmt-layer-adaptive-rpi5}.

\bmhead{Code availability}
The validated heterogeneous XMSSMT source modifications, analysis scripts, provenance records, and reproducibility artifacts supporting this study are available in the public repository at \url{https://github.com/sohaibrana1/xmssmt-layer-adaptive-rpi5}.

\bmhead{Competing interests}
The authors declare that they have no competing interests.

\bibliography{references}

\end{document}
